스키마에 부합하는 게임 내 행동 제안이라도 접근할 수 없는 물체를 요구하거나 권한이 없는 상태 전이, 혹은 현재 시점에서 공개할 수 없는 비밀 발화를 포함할 수 있다. 서사적 정합성(grounding)과 개연성 있는 대화 생성 \cite{S03_weir2024ontologically,S20_ashby2023personalized}은 후보 콘텐츠를 제공할 뿐이며, 게임 엔진은 현재 상태에서 어떤 전이를 승인할지 별도로 판정해야 한다.

\system{}는 이 판정을 외부의 제안--커밋(proposal-to-commit) 경계에서 강제한다(그림~\ref{fig:architecture}). 모델과 어댑터가 구조화된 후보를 제출하면 결정론적 게이트가 행동 정책, 선행조건, 물체 도달 가능성, NPC 지식, 정보 공개 권한, 효과, 퀘스트 단계 단조성을 검사한다. 검사에 실패하면 게이트는 제한된 수리를 위한 구조화된 오류를 반환하며, 수리는 정본 상태(canonical state)가 아닌 후보만 수정하고 동일 상태 재검증을 통과한 경우에만 상태를 갱신한다. 여덟 역할 읽기 전용 온톨로지 뷰는 쓰기 권한 없이 출처 선언과 검사·결정을 연결하며, 해시 연결 전이 기록은 의미적 리플레이와 배정 전수 회계(assignment-complete accounting)를 보장한다.

개념적으로 이 계약은 단순 사실 일치(factual grounding)를 넘어 NPC의 \emph{인식론적 지식 경계}(행위자가 알고 있는 사실)와 \emph{정보 공개 권한 정책}(현재 퀘스트 단계에서 발화할 수 있는 사실)을 분리한다. 방법론적으로는 LLM 호출이 0회인 결정론적 검증 계층(게이트 검사, 유도 수리, 엔진 리플레이)과 실패를 누락 없이 보존하는 제한된 확률적 제안기 계층을 엄격히 분리한다.

결정론적 감사에서 스키마 전용 대조군은 32개 유한 그래프 조합에서 무효 제안 260/416건을 커밋한 반면 의미 게이트는 \textbf{0/416} 위험 커밋·0/156 오거부·128/260건 $\rho$ 수리를 기록했고, $R{\le}64, d{\le}5$ 480건 Datalog 검증(\textbf{P9})에서 \texttt{nesy\_game.multihop}만이 $\le$7.12\,$\mu$s로 100.0\% 정확도를 달성하였다($Q{=}1100.19$). 정리 1--4(\textbf{I1--I4})를 토대로, Haiku·Sonnet·Opus~5.5 직접 HTTPS Anthropic API 평가는 105회 독립 5턴 폐루프 협상(\textbf{P15}: 15/15 무누출·0 위험·하이브리드 15/15 목표 대 Safe-RL~\cite{alshiekh2018shield} 1/15·NeMo 스타일~\cite{rebedea2023nemo} 1/15), 60회 요인 실험(\textbf{P16A}: \texttt{ontology\_multihop} 30/30 무누출·22건 수리 대 22/30 \texttt{general\_rule}), 21회 Godot~4.7.1 폐루프(\textbf{P16B}: 21/21 재실행 일치, 제품 6/6 해시 일치), 540회 통합 한·영 우회(\textbf{P17}: $\kappa{=}0.8813$, 0\% FRR)를 입증하며, 3모델군 CLI(\textbf{P11}--\textbf{P14}: $(a,E)$ 수리 27/27 0토큰), \texttt{DATASHEET\_v5.md}~\cite{gebru2021datasheets}, Croissant~1.0 v5~\cite{akhtar2024croissant}(26/26), 격리된 \texttt{TR01--TR06}($N{=}0$, \texttt{AWAITING\_HUMAN\_RATERS})으로 공개 경계를 완결한다.

\section{관련 연구와 연구 공백}
정합성, 기억, 계획, 역할 기반 에이전트 \cite{S03_weir2024ontologically,S20_ashby2023personalized,S07_park2023generative,S08_shao2023character} 및 플레이어·대화 시스템 \cite{S21_akoury2023player,S22_hochreiter2026beyond,S23_yin2026contextualized,S02_figueiredo2025scaffolded}은 맥락적 상호작용을 지원하는 반면, 본 연구는 인코딩된 현재 상태와 공개 정책이 제안된 전이를 허용하는지를 판정한다. 정서 및 몰입 모델링 \cite{S24_yannakakis2011experiencedriven,S25_melhart2020engagement,S26_wang2026engagement}은 생성 스타일을 안내할 뿐 상태 변경 권한을 부여하지 않는다.

기존 벤치마크와 자동 플레이테스팅은 상태, 체화 환경, 콘텐츠, 과정, 사회적 상호작용, 장기 과업을 다루며 \cite{S11_cote2019textworld,S12_wang2022scienceworld,S13_fan2022minedojo,S14_perezliebana2019gvgai,S15_liu2024agentbench,S16_ma2024agentboard,S17_zhou2024sotopia,S18_paglieri2025balrog,S19_lepelletier2022playtesting}, NCP-Bench는 서사적 일관성과 사실 정합성을 평가한다 \cite{S58_ma2026ncpbench}. 문법·스키마·어포던스·안전 실딩 기법 \cite{S09_geng2023grammar,S10_beurerkellner2023prompting,S36_scholak2021picard,S50_kambhampati2024llmmodulo,S51_ichter2023saycan,S52_wang2025agentspec,S53_li2025settingdc,S54_zhou2025story2game,S42_buongiorno2024pangea,alshiekh2018shield,rebedea2023nemo}, 반례 기반 수리 \cite{S45_madaan2023selfrefine,S48_shinn2023reflexion,S46_chen2023selfdebug,S47_solarlezama2006sketching}, 재현성·데이터셋 표준 \cite{S31_agarwal2021precipice,S33_pineau2021reproducibility,S56_feng2025agentrr,S59_sun2025bdc,gebru2021datasheets,akhtar2024croissant}이 본 설계의 토대다. 스키마 적합성과 행위 합법성, 또는 사실 정합성과 발화 안전성을 동일시하던 선행 연구와 달리, \system{}는 스키마 파싱, 인식론적 지식, 공개 권한, 정본 상태 커밋을 독립 감사 가능한 게이트로 분리한다.
